% paper10_fermion_mass_sketch.tex
% Fermion Mass as Retrodiction Cost in Internal Space
% Author: Sheng-Kai Huang
% Compile: pdflatex paper10_fermion_mass_sketch.tex
% Target: 3 pages sketch, PRL style

\documentclass[prl,twocolumn,superscriptaddress,longbibliography]{revtex4-2}

\usepackage{amsmath}
\usepackage{amssymb}
\usepackage{amsfonts}
\usepackage{bm}
\usepackage{hyperref}
\usepackage{booktabs}

% Macros
\newcommand{\kB}{k_{\mathrm{B}}}
\newcommand{\Tr}{\mathrm{Tr}}
\newcommand{\tr}{\mathrm{tr}}
\newcommand{\SvN}{S_{\mathrm{vN}}}
\newcommand{\AF}{A_{\mathrm{F}}}
\newcommand{\DF}{D_{\mathrm{F}}}
\newcommand{\HF}{H_{\mathrm{F}}}
\newcommand{\Mref}{M_{\mathrm{ref}}}
\newcommand{\JO}{J_3(\mathbb{O})}
\newcommand{\Petz}{\mathcal{R}}

\begin{document}

\title{Fermion Mass as Retrodiction Cost:\\
  $\Sigma$ on $J_3(\mathbb{O})$ and the Electroweak Angle}

\author{Sheng-Kai Huang}
\email{akai@fawstudio.com}
\affiliation{Independent Researcher}

\date{\today}

\begin{abstract}
We propose that fermion mass is the cost of retrodiction
failure in the internal (gauge) space of the Standard Model.
In the $\Sigma = 2\ln Q$ framework, the temporal asymmetry
$\tau_f = 1 - F_f$ of fermion~$f$ in the internal spectral
triple $(\AF, \HF, \DF)$ determines its mass via
$m_f = \Mref\,\tau_f$.
We identify the internal quantum relative entropy
$\Sigma_f = D(\rho_f \| \rho_0)$ with the eigenvalue spectrum
of the exceptional Jordan algebra~$\JO$, connecting to Singh's
parameter-free mass ratios~[CONJECTURED].
A key observation: Singh's universal parameter $\delta^2 = 3/8$
for $\JO$ eigenvalue spreading \emph{numerically equals}
the Weinberg angle $\sin^2\theta_W = 3/8$ at GUT scale
in the Chamseddine--Connes noncommutative geometry~[CONJECTURED].
If these are algebraically identical---both arising from the
maximal associative subalgebra $\AF \subset \JO$
(Boyle~2020)---then fermion mass ratios and the electroweak
mixing angle share a common information-theoretic origin.
We outline the derivation program and state clearly which
steps are proven, conjectured, or speculative.
\end{abstract}

\maketitle


%=======================================================================
\section{Motivation: Retrodiction cost in internal space}
\label{sec:motivation}
%=======================================================================

In companion papers~\cite{PaperI,PaperII,Paper9}, we established
that the quantum relative entropy (QRE)
\begin{equation}\label{eq:Sigma}
  \Sigma = D\bigl(\rho_{\mathrm{st}}
  \,\big\|\, \rho_{\mathrm{mat}}\bigr) \geq 0
\end{equation}
governs the arrow of time ($\tau = 1 - F$, Paper~I~[PROVEN]),
the gravitational metric ($\Sigma_{\mathrm{grav}} = 2\ln Q$,
Paper~II~[PROVEN]), and the dimensionality of spacetime
($d=4$ from the spectral Fisher metric,
Paper~9~[PROVEN]).

The Standard Model lives on the
almost-commutative geometry $M \times F$ of
Connes~\cite{CC1996}, where $F$ is the finite
internal space with algebra
$\AF = \mathbb{C} \oplus \mathbb{H}
\oplus M_3(\mathbb{C})$.
The finite Dirac operator~$\DF$ encodes Yukawa
couplings, hence fermion masses~[KNOWN].

We now extend $\Sigma$ to the internal space:
\begin{equation}\label{eq:Sigma_f}
  \Sigma_f \;=\; D\bigl(\rho_f \,\big\|\, \rho_0\bigr)_F\,,
\end{equation}
where $\rho_f$ is the state of fermion~$f$ in~$\HF$,
$\rho_0 = I/\dim(\HF)$ is the vacuum reference, and
the subscript~$F$ denotes restriction to the internal
algebra~[CONJECTURED].

The fermion mass is then the
retrodiction cost:
\begin{equation}\label{eq:mass_formula}
  \boxed{\;
  m_f \;=\; \Mref \times \tau_f
  \;=\; \Mref \times
  \bigl(1 - e^{-\Sigma_f/2}\bigr)
  \;}
\end{equation}
with $\Mref$ a reference scale identified with the
electroweak VEV $v = 246$\,GeV~[CONJECTURED].

Eq.~\eqref{eq:mass_formula} has the correct limiting
behavior:
\begin{itemize}
\item $\Sigma_f = 0$: massless (perfect retrodiction,
  $F_f = 1$, e.g., unbroken gauge symmetry);
\item $\Sigma_f \gg 1$: $m_f \to \Mref$ (saturation;
  the top quark at $m_t/v \approx 0.70$ is near this
  limit);
\item $\Sigma_f \ll 1$: $m_f \approx \Mref\,\Sigma_f/2$
  (linear regime for light fermions).
\end{itemize}


%=======================================================================
\section{$\Sigma_f$ from the spectral action}
\label{sec:spectral}
%=======================================================================

Chamseddine, Connes, and van
Suijlekom~\cite{CCSvS2018} proved that the von~Neumann
entropy of the fermionic second quantization of a
spectral triple equals a spectral action:
$\SvN(\rho_\beta) = \Tr(h(\beta D))$~[KNOWN].

The spectral action on $F$ expands as
\begin{equation}\label{eq:SF}
  S_F = \Tr\bigl(f(\DF^2/\Lambda^2)\bigr)
  \;\approx\; f_0\,\Tr(1)
  + f_2\,\Lambda^{-2}\,\Tr(\DF^2) + \cdots
\end{equation}
The leading mass-dependent term is~[KNOWN]
\begin{equation}\label{eq:TrDF2}
  \Tr(\DF^2) = \sum_f N_f\, m_f^2
  = v^2 \sum_f N_f\, y_f^2\,,
\end{equation}
where $N_f$ counts color and chirality
multiplicities and $y_f$ is the Yukawa coupling.

Identifying $\Sigma_f$ with the contribution of
fermion~$f$ to the QRE of the internal state:
\begin{equation}\label{eq:Sigma_Yukawa}
  \Sigma_f \;\propto\; y_f^2
  \;=\; \bigl(m_f/v\bigr)^2
  \qquad\text{[CONJECTURED]}
\end{equation}
This is consistent with Eq.~\eqref{eq:mass_formula}
in the light-fermion regime ($\Sigma_f \ll 1$):
$m_f \approx v\,\Sigma_f/2 \propto v\, y_f^2/2$,
which gives $m_f \propto y_f^2$---differing from
the usual $m_f = y_f v/\sqrt{2}$.
The resolution is that
Eq.~\eqref{eq:mass_formula} determines $\tau_f$
(the dimensionless asymmetry), not $m_f$ directly.
The correct identification is:
\begin{equation}\label{eq:tau_y}
  \tau_f = y_f / \sqrt{2}
  \quad\Longrightarrow\quad
  \Sigma_f = -2\ln(1 - y_f/\sqrt{2})\,.
\end{equation}
For the top quark, $y_t/\sqrt{2} \approx 0.70$,
giving $\Sigma_t \approx 2.41$---a large but finite
retrodiction cost~[CONJECTURED].


%=======================================================================
\section{Connection to $J_3(\mathbb{O})$ mass ratios}
\label{sec:J3O}
%=======================================================================

Singh~\cite{Singh2025} showed that the complexified
exceptional Jordan algebra $J_3(\mathbb{O}_\mathbb{C})$
provides parameter-free fermion mass ratios with a
single universal parameter $\delta^2 = 3/8$.
Three generations arise as three eigenvalues of a Jordan
element $X \in \JO$; the mass hierarchy comes from a
``minimal universal ladder'' in $\mathrm{Sym}^3(3)$
of the flavor $\mathrm{SU}(3)$~[KNOWN].

We observe that $\delta^2 = 3/8$ \emph{numerically
equals} the Weinberg angle at GUT scale in the
Chamseddine--Connes framework~\cite{CC1996}:
\begin{equation}\label{eq:38}
  \sin^2\theta_W\big|_{\mathrm{GUT}}
  = \frac{\Tr(T_3^2)}{\Tr(T_3^2) + \Tr(Y^2)}
  = \frac{3}{8}\,.
\end{equation}

Both arise from the algebraic structure of the
Standard Model:
\begin{itemize}
\item $\sin^2\theta_W = 3/8$: from the fermion
  hypercharge assignments dictated by
  $\AF = \mathbb{C} \oplus \mathbb{H}
  \oplus M_3(\mathbb{C})$~[KNOWN].
\item $\delta^2 = 3/8$: from the
  $\mathrm{Sym}^3(3)$ structure in
  $J_3(\mathbb{O}_\mathbb{C})$~[KNOWN,
  Singh~2025].
\end{itemize}

Boyle~\cite{Boyle2020} showed that
$\AF$ is the maximal associative subalgebra of~$\JO$.
This establishes that both $3/8$'s live in the
\emph{same algebraic structure}~[KNOWN].

\medskip
\noindent\textbf{Conjecture.}\;
\emph{The two instances of $3/8$ are algebraically
identical: $\delta^2 = \sin^2\theta_W$ because both
express the ratio of abelian to total gauge content
within $\AF \subset \JO$.}
~[CONJECTURED]

If true, the electroweak mixing angle and the
fermion mass hierarchy share a common origin in
the exceptional Jordan algebra---unified by the
retrodiction cost $\Sigma$ on the internal space.


%=======================================================================
\section{The Petz map on $J_3(\mathbb{O})$}
\label{sec:petz_J3O}
%=======================================================================

The Petz recovery map $\Petz$ is the unique
retrodiction functor satisfying the detailed balance
condition~\cite{PaperI,Petz1986}.
On $\JO$, we define~[SPECULATIVE]:
\begin{equation}
  \Petz_\sigma \;\colon\;
  J_3(\mathbb{O}) \to J_3(\mathbb{O})\,,
  \qquad
  \Petz_\sigma(X) = \sigma^{1/2}\,X\,\sigma^{1/2}\,,
\end{equation}
where $\sigma$ is a reference state in the Jordan
algebra (the product is the Jordan product
$A \circ B = \tfrac{1}{2}(AB + BA)$).

The recovery fidelity for generation~$i$:
\begin{equation}
  F_i = F\bigl(\rho_i,\, \Petz_\sigma(\rho_i)\bigr)
  = e^{-\Sigma_i/2}\,.
\end{equation}

\noindent\textbf{Speculative identification.}\;
Singh's ``minimal universal ladder'' = the Petz
recovery channel $\Petz_\sigma$ on $\JO$ that
maximizes the average fidelity
$\bar{F} = \frac{1}{3}\sum_{i=1}^{3} F_i$
subject to the Jordan algebra constraints.

If this holds, the mass ratios follow from Petz
optimality, and the CKM matrix elements are
the fidelity overlaps between up-type and
down-type Petz channels:
\begin{equation}
  |V_{ij}|^2 = F\bigl(\Petz_\sigma^{(u,i)},\,
  \Petz_\sigma^{(d,j)}\bigr)
  \qquad\text{[SPECULATIVE]}\,.
\end{equation}


%=======================================================================
\section{Mass table and predictions}
\label{sec:predictions}
%=======================================================================

Table~\ref{tab:masses} lists the retrodiction
cost~$\Sigma_f$ for each SM fermion, computed from
Eq.~\eqref{eq:tau_y} with $\tau_f = m_f/v$.

\begin{table}[t]
\centering
\caption{Retrodiction cost $\Sigma_f$ and temporal
  asymmetry $\tau_f$ for SM fermions, using
  $\Mref = v = 246$\,GeV.
  Values computed from $\tau_f = m_f/v$ and
  $\Sigma_f = -2\ln(1 - \tau_f)$.}
\label{tab:masses}
\begin{ruledtabular}
\begin{tabular}{lccc}
Fermion & $m_f$ (GeV) & $\tau_f$ & $\Sigma_f$ \\
\hline
$\nu$ (lightest) & $\lesssim 10^{-11}$ & $\lesssim 4\times 10^{-14}$ & $\lesssim 10^{-13}$ \\
$e$   & $5.11\times10^{-4}$ & $2.08\times10^{-6}$ & $4.15\times10^{-6}$ \\
$u$   & $2.2\times10^{-3}$  & $8.9\times10^{-6}$  & $1.79\times10^{-5}$ \\
$d$   & $4.7\times10^{-3}$  & $1.9\times10^{-5}$  & $3.83\times10^{-5}$ \\
$\mu$ & $0.106$             & $4.31\times10^{-4}$  & $8.62\times10^{-4}$ \\
$s$   & $0.093$             & $3.78\times10^{-4}$  & $7.57\times10^{-4}$ \\
$c$   & $1.27$              & $5.16\times10^{-3}$  & $1.03\times10^{-2}$ \\
$\tau$& $1.777$             & $7.22\times10^{-3}$  & $1.45\times10^{-2}$ \\
$b$   & $4.18$              & $1.70\times10^{-2}$  & $3.43\times10^{-2}$ \\
$t$   & $173$               & $0.703$              & $2.41$ \\
\end{tabular}
\end{ruledtabular}
\end{table}

\noindent The top quark has $\Sigma_t = 2.41$, making it
the most ``internally irreversible'' fermion.
The electron has $\Sigma_e = 4.15 \times 10^{-6}$: it
is nearly perfectly retrodictable in the internal space.
The neutrino has $\Sigma_\nu \lesssim 10^{-13}$: it
is the closest to a ``zero-entropy'' internal state.

\medskip
\noindent\textbf{Testable prediction.}\;
$m_f \leq v$ for all SM fermions (Petz bound saturation
with $\Mref = v$). This is satisfied: $m_t = 173\,\mathrm{GeV}
< 246\,\mathrm{GeV}$~[CONJECTURED, consistent with data].

\medskip
\noindent\textbf{Open prediction.}\;
If the $\delta^2 = \sin^2\theta_W$ identification holds,
and Singh's mass ratios are correct, then the $\Sigma$
framework \emph{with no free parameters} predicts all
fermion mass ratios from the algebraic structure of
$\JO$.


%=======================================================================
\section{Discussion and status}
\label{sec:discussion}
%=======================================================================

We summarize the logical status of each claim:

\begin{description}
\item[PROVEN:]
  $\tau = 1-F$, $\Sigma_{\mathrm{grav}} = 2\ln Q$,
  $g_{QQ} = 2/Q^2$ in $d=4$,
  $\sin^2\theta_W = 3/8$ in NCG,
  spectral action = $\SvN$,
  Singh's mass ratios from $\JO$.

\item[CONJECTURED:]
  $\Sigma_f = D(\rho_f\|\rho_0)_F$ defines the
  internal retrodiction cost;
  $m_f = v\,\tau_f$;
  $\delta^2 = \sin^2\theta_W$ is algebraically the same
  $3/8$;
  the Petz bound gives $m_f \leq v$.

\item[SPECULATIVE:]
  $\Sigma$ extremization on $\JO$ selects the physical
  mass ratios;
  Singh's ladder = Petz recovery on $\JO$;
  CKM matrix = Petz fidelity between up/down channels.
\end{description}

The key testable step is the algebraic identification
of the two $3/8$'s.
If the Boyle embedding $\AF \subset \JO$
maps Singh's $\delta^2$ to the NCG $\sin^2\theta_W$,
this would unify the electroweak mixing angle and the
fermion mass hierarchy within a single
information-theoretic framework.

Concretely, one must show that the
$\mathrm{Sym}^3(3)$ eigenvalue spread parameter of
$\JO$ reduces to $\Tr(T_3^2)/[\Tr(T_3^2)+\Tr(Y^2)]$
when restricted to $\AF$.
This is a finite-dimensional algebraic computation
that we leave to a subsequent paper.

\begin{acknowledgments}
The author thanks the open-source community for the
division algebra and NCG literature, and acknowledges
the foundational work of Connes, Singh, Furey, Boyle,
and Farnsworth.
\end{acknowledgments}


\begin{thebibliography}{20}

\bibitem{PaperI}
S.-K.~Huang, ``Temporal asymmetry as Petz recovery
fidelity,'' (Paper~I).

\bibitem{PaperII}
S.-K.~Huang, ``The gravitational refractive index as
Petz recovery fidelity,'' (Paper~II).

\bibitem{Paper9}
S.-K.~Huang, ``Temporal asymmetry as spectral Fisher
information: Why four dimensions?'' (Paper~9).

\bibitem{CC1996}
A.~H.~Chamseddine and A.~Connes,
``The spectral action principle,''
Commun.\ Math.\ Phys.\  \textbf{186}, 731 (1997),
\href{https://arxiv.org/abs/hep-th/9606001}{hep-th/9606001}.

\bibitem{CCSvS2018}
A.~H.~Chamseddine, A.~Connes, and
W.~D.~van~Suijlekom,
``Entropy and the spectral action,''
Commun.\ Math.\ Phys.\  \textbf{373}, 689 (2020),
\href{https://arxiv.org/abs/1809.02944}{arXiv:1809.02944}.

\bibitem{Singh2025}
T.~P.~Singh,
``Fermion mass ratios from the exceptional Jordan
algebra,''
\href{https://arxiv.org/abs/2508.10131}{arXiv:2508.10131}
(2025).

\bibitem{Boyle2020}
L.~Boyle,
``The Standard Model, the exceptional Jordan algebra,
and triality,''
\href{https://arxiv.org/abs/2006.16265}{arXiv:2006.16265}
(2020).

\bibitem{Furey2024}
C.~Furey and M.~J.~Hughes,
``Three generations and a trio of trialities,''
Phys.\ Lett.\ B (2025),
\href{https://arxiv.org/abs/2409.17948}{arXiv:2409.17948}.

\bibitem{Farnsworth2026}
S.~Farnsworth, F.~Finster, C.~Paganini, and
T.~P.~Singh,
``Causal fermion systems, non-commutative geometry and
generalized trace dynamics,''
\href{https://arxiv.org/abs/2603.05018}{arXiv:2603.05018}
(2026).

\bibitem{DorauMuch2025}
A.~Dorau and A.~Much,
``From quantum relative entropy to the semiclassical
Einstein equations,''
Phys.\ Rev.\ Lett.\  (2025),
\href{https://arxiv.org/abs/2510.24491}{arXiv:2510.24491}.

\bibitem{Petz1986}
D.~Petz,
``Sufficient subalgebras and the relative entropy of
states of a von Neumann algebra,''
Commun.\ Math.\ Phys.\  \textbf{105}, 123 (1986).

\bibitem{Szangolies2025}
J.~Szangolies,
``From qubits to the Standard Model,''
\href{https://arxiv.org/abs/2512.17328}{arXiv:2512.17328}
(2025).

\bibitem{Kaplan1992}
D.~B.~Kaplan,
``A method for simulating chiral fermions on the
lattice,''
Phys.\ Lett.\ B \textbf{288}, 342 (1992).

\bibitem{Furey2025super}
C.~Furey,
``A superalgebra within,''
Annalen der Physik (2025),
\href{https://arxiv.org/abs/2505.07923}{arXiv:2505.07923}.

\end{thebibliography}

\end{document}
